% Lecture 10 - Page 9 (6-XII-1994).
% Fragment for % Lecture 10 - Page 9 (6-XII-1994).
% Fragment for % Lecture 10 - Page 9 (6-XII-1994).
% Fragment for % Lecture 10 - Page 9 (6-XII-1994).
% Fragment for \input{courses/pages/curs10-pagina9-en}.
% No preamble and no forced page break.
% Continuation of the proof of Sylow's First Theorem from Page 8.
% Editorial notes flag doubtful readings rather than silently correcting them.

\subsection*{Lecture 10 - Page 9}

\noindent $\varphi(K)=\dfrac{K}{H}=K'.$

\medskip
\noindent Also $H\trianglelefteq K$, since $H\trianglelefteq N_G(H)$.

\medskip
\noindent Apply Lagrange's theorem:
\[
|K|=|H|\cdot[K:H]=|H|\cdot[\,\ldots\,]
\quad\Longrightarrow\quad p\mid|K'|.
\]
% In the manuscript, the final index in square brackets is omitted/unclear.
% The last symbols appear to read "p | |K'|"; mathematically the aim
% is to establish |K|=p^{i+1}. The apparent reading has not been replaced.

\bigskip
\noindent\fbox{\textbf{Proof of Sylow's Second Theorem}}

\medskip
\noindent Let $P\in\operatorname{Syl}_p(G)$. Let
\[
A=\{aP\mid a\in G\}.
\]

\noindent $H\times A\longrightarrow A$,
\[
(h,aP)\longmapsto(ha)P\in A,
\]
\noindent an action by translation.

\noindent Verify that this is an action.

\medskip
\noindent Consider the corresponding $A_0$:
\[
A_0=\{aP\mid haP=aP,\ (\forall)\,h\in H\}.
\]

\noindent $A_0\neq\varnothing$, since $p\nmid[G:P]$ and
$|P|=p^t$.

\medskip
\noindent By Lemma 2 $\Longrightarrow$
\[
p\nmid[N_G(P):P].
\]
% The text appears to show "[N_G(P):P]", but when Lemma 2 is
% applied to H acting on G/P, the immediate conclusion concerns |A_0|.
% This apparent inconsistency is retained for manuscript review.

% The line "Apply Lemma 2" is crossed out in the manuscript.
\medskip
\noindent
\[
|A_0|\equiv|A|\pmod p
\quad\Longrightarrow\quad p\nmid|A_0|
\quad\Longrightarrow\quad A_0\neq\varnothing.
\]
\noindent $|A|=m$ and $p\nmid m$.

\medskip
\noindent There exists $aP\in A_0$. Then
\[
haP=aP,\quad(\forall)\,h\in H
\quad\Longleftrightarrow\quad
 a^{-1}haP=P,\quad(\forall)\,h\in H.
\]
\[
\Longleftrightarrow\quad a^{-1}Ha\leq P
\quad\Longrightarrow\quad H\leq aPa^{-1}.
\]

\medskip
\noindent Let $Q\in\operatorname{Syl}_p(G)$. Then
\[
aQa^{-1}\leq P.
\]
\noindent But
\[
|P|=p^t,\qquad |Q|=p^t
\quad\Longrightarrow\quad |aQa^{-1}|=p^t
\quad\Longrightarrow\quad aQa^{-1}=P.
\]

\bigskip
\noindent\underline{Application}: $|G|=pq$, $p>q$, $q\nmid(p-1)$.

\medskip
\noindent By Cauchy's theorem, there exist $a,b\in G$ with
$o(a)=p$, $o(b)=q$.

\noindent $P=\langle a\rangle$, $Q=\langle b\rangle$.

\noindent By Sylow's Third Theorem:
\[
|\operatorname{Syl}_p(G)|\mid pq,
\qquad |\operatorname{Syl}_p(G)|\equiv1\pmod p
\quad\Longrightarrow\quad |\operatorname{Syl}_p(G)|=1.
\]

\noindent $P\in\operatorname{Syl}_p(G)\ \ni Q$.%
% The final reading "\ni Q" is uncertain. If interpreted as
% containment Q<=P, it conflicts with |Q|=q and |P|=p.
% The manuscript page ends here; no additional proof is supplied.
.
% No preamble and no forced page break.
% Continuation of the proof of Sylow's First Theorem from Page 8.
% Editorial notes flag doubtful readings rather than silently correcting them.

\subsection*{Lecture 10 - Page 9}

\noindent $\varphi(K)=\dfrac{K}{H}=K'.$

\medskip
\noindent Also $H\trianglelefteq K$, since $H\trianglelefteq N_G(H)$.

\medskip
\noindent Apply Lagrange's theorem:
\[
|K|=|H|\cdot[K:H]=|H|\cdot[\,\ldots\,]
\quad\Longrightarrow\quad p\mid|K'|.
\]
% In the manuscript, the final index in square brackets is omitted/unclear.
% The last symbols appear to read "p | |K'|"; mathematically the aim
% is to establish |K|=p^{i+1}. The apparent reading has not been replaced.

\bigskip
\noindent\fbox{\textbf{Proof of Sylow's Second Theorem}}

\medskip
\noindent Let $P\in\operatorname{Syl}_p(G)$. Let
\[
A=\{aP\mid a\in G\}.
\]

\noindent $H\times A\longrightarrow A$,
\[
(h,aP)\longmapsto(ha)P\in A,
\]
\noindent an action by translation.

\noindent Verify that this is an action.

\medskip
\noindent Consider the corresponding $A_0$:
\[
A_0=\{aP\mid haP=aP,\ (\forall)\,h\in H\}.
\]

\noindent $A_0\neq\varnothing$, since $p\nmid[G:P]$ and
$|P|=p^t$.

\medskip
\noindent By Lemma 2 $\Longrightarrow$
\[
p\nmid[N_G(P):P].
\]
% The text appears to show "[N_G(P):P]", but when Lemma 2 is
% applied to H acting on G/P, the immediate conclusion concerns |A_0|.
% This apparent inconsistency is retained for manuscript review.

% The line "Apply Lemma 2" is crossed out in the manuscript.
\medskip
\noindent
\[
|A_0|\equiv|A|\pmod p
\quad\Longrightarrow\quad p\nmid|A_0|
\quad\Longrightarrow\quad A_0\neq\varnothing.
\]
\noindent $|A|=m$ and $p\nmid m$.

\medskip
\noindent There exists $aP\in A_0$. Then
\[
haP=aP,\quad(\forall)\,h\in H
\quad\Longleftrightarrow\quad
 a^{-1}haP=P,\quad(\forall)\,h\in H.
\]
\[
\Longleftrightarrow\quad a^{-1}Ha\leq P
\quad\Longrightarrow\quad H\leq aPa^{-1}.
\]

\medskip
\noindent Let $Q\in\operatorname{Syl}_p(G)$. Then
\[
aQa^{-1}\leq P.
\]
\noindent But
\[
|P|=p^t,\qquad |Q|=p^t
\quad\Longrightarrow\quad |aQa^{-1}|=p^t
\quad\Longrightarrow\quad aQa^{-1}=P.
\]

\bigskip
\noindent\underline{Application}: $|G|=pq$, $p>q$, $q\nmid(p-1)$.

\medskip
\noindent By Cauchy's theorem, there exist $a,b\in G$ with
$o(a)=p$, $o(b)=q$.

\noindent $P=\langle a\rangle$, $Q=\langle b\rangle$.

\noindent By Sylow's Third Theorem:
\[
|\operatorname{Syl}_p(G)|\mid pq,
\qquad |\operatorname{Syl}_p(G)|\equiv1\pmod p
\quad\Longrightarrow\quad |\operatorname{Syl}_p(G)|=1.
\]

\noindent $P\in\operatorname{Syl}_p(G)\ \ni Q$.%
% The final reading "\ni Q" is uncertain. If interpreted as
% containment Q<=P, it conflicts with |Q|=q and |P|=p.
% The manuscript page ends here; no additional proof is supplied.
.
% No preamble and no forced page break.
% Continuation of the proof of Sylow's First Theorem from Page 8.
% Editorial notes flag doubtful readings rather than silently correcting them.

\subsection*{Lecture 10 - Page 9}

\noindent $\varphi(K)=\dfrac{K}{H}=K'.$

\medskip
\noindent Also $H\trianglelefteq K$, since $H\trianglelefteq N_G(H)$.

\medskip
\noindent Apply Lagrange's theorem:
\[
|K|=|H|\cdot[K:H]=|H|\cdot[\,\ldots\,]
\quad\Longrightarrow\quad p\mid|K'|.
\]
% In the manuscript, the final index in square brackets is omitted/unclear.
% The last symbols appear to read "p | |K'|"; mathematically the aim
% is to establish |K|=p^{i+1}. The apparent reading has not been replaced.

\bigskip
\noindent\fbox{\textbf{Proof of Sylow's Second Theorem}}

\medskip
\noindent Let $P\in\operatorname{Syl}_p(G)$. Let
\[
A=\{aP\mid a\in G\}.
\]

\noindent $H\times A\longrightarrow A$,
\[
(h,aP)\longmapsto(ha)P\in A,
\]
\noindent an action by translation.

\noindent Verify that this is an action.

\medskip
\noindent Consider the corresponding $A_0$:
\[
A_0=\{aP\mid haP=aP,\ (\forall)\,h\in H\}.
\]

\noindent $A_0\neq\varnothing$, since $p\nmid[G:P]$ and
$|P|=p^t$.

\medskip
\noindent By Lemma 2 $\Longrightarrow$
\[
p\nmid[N_G(P):P].
\]
% The text appears to show "[N_G(P):P]", but when Lemma 2 is
% applied to H acting on G/P, the immediate conclusion concerns |A_0|.
% This apparent inconsistency is retained for manuscript review.

% The line "Apply Lemma 2" is crossed out in the manuscript.
\medskip
\noindent
\[
|A_0|\equiv|A|\pmod p
\quad\Longrightarrow\quad p\nmid|A_0|
\quad\Longrightarrow\quad A_0\neq\varnothing.
\]
\noindent $|A|=m$ and $p\nmid m$.

\medskip
\noindent There exists $aP\in A_0$. Then
\[
haP=aP,\quad(\forall)\,h\in H
\quad\Longleftrightarrow\quad
 a^{-1}haP=P,\quad(\forall)\,h\in H.
\]
\[
\Longleftrightarrow\quad a^{-1}Ha\leq P
\quad\Longrightarrow\quad H\leq aPa^{-1}.
\]

\medskip
\noindent Let $Q\in\operatorname{Syl}_p(G)$. Then
\[
aQa^{-1}\leq P.
\]
\noindent But
\[
|P|=p^t,\qquad |Q|=p^t
\quad\Longrightarrow\quad |aQa^{-1}|=p^t
\quad\Longrightarrow\quad aQa^{-1}=P.
\]

\bigskip
\noindent\underline{Application}: $|G|=pq$, $p>q$, $q\nmid(p-1)$.

\medskip
\noindent By Cauchy's theorem, there exist $a,b\in G$ with
$o(a)=p$, $o(b)=q$.

\noindent $P=\langle a\rangle$, $Q=\langle b\rangle$.

\noindent By Sylow's Third Theorem:
\[
|\operatorname{Syl}_p(G)|\mid pq,
\qquad |\operatorname{Syl}_p(G)|\equiv1\pmod p
\quad\Longrightarrow\quad |\operatorname{Syl}_p(G)|=1.
\]

\noindent $P\in\operatorname{Syl}_p(G)\ \ni Q$.%
% The final reading "\ni Q" is uncertain. If interpreted as
% containment Q<=P, it conflicts with |Q|=q and |P|=p.
% The manuscript page ends here; no additional proof is supplied.
.
% No preamble and no forced page break.
% Continuation of the proof of Sylow's First Theorem from Page 8.
% Editorial notes flag doubtful readings rather than silently correcting them.

\subsection*{Lecture 10 - Page 9}

\noindent $\varphi(K)=\dfrac{K}{H}=K'.$

\medskip
\noindent Also $H\trianglelefteq K$, since $H\trianglelefteq N_G(H)$.

\medskip
\noindent Apply Lagrange's theorem:
\[
|K|=|H|\cdot[K:H]=|H|\cdot[\,\ldots\,]
\quad\Longrightarrow\quad p\mid|K'|.
\]
% In the manuscript, the final index in square brackets is omitted/unclear.
% The last symbols appear to read "p | |K'|"; mathematically the aim
% is to establish |K|=p^{i+1}. The apparent reading has not been replaced.

\bigskip
\noindent\fbox{\textbf{Proof of Sylow's Second Theorem}}

\medskip
\noindent Let $P\in\operatorname{Syl}_p(G)$. Let
\[
A=\{aP\mid a\in G\}.
\]

\noindent $H\times A\longrightarrow A$,
\[
(h,aP)\longmapsto(ha)P\in A,
\]
\noindent an action by translation.

\noindent Verify that this is an action.

\medskip
\noindent Consider the corresponding $A_0$:
\[
A_0=\{aP\mid haP=aP,\ (\forall)\,h\in H\}.
\]

\noindent $A_0\neq\varnothing$, since $p\nmid[G:P]$ and
$|P|=p^t$.

\medskip
\noindent By Lemma 2 $\Longrightarrow$
\[
p\nmid[N_G(P):P].
\]
% The text appears to show "[N_G(P):P]", but when Lemma 2 is
% applied to H acting on G/P, the immediate conclusion concerns |A_0|.
% This apparent inconsistency is retained for manuscript review.

% The line "Apply Lemma 2" is crossed out in the manuscript.
\medskip
\noindent
\[
|A_0|\equiv|A|\pmod p
\quad\Longrightarrow\quad p\nmid|A_0|
\quad\Longrightarrow\quad A_0\neq\varnothing.
\]
\noindent $|A|=m$ and $p\nmid m$.

\medskip
\noindent There exists $aP\in A_0$. Then
\[
haP=aP,\quad(\forall)\,h\in H
\quad\Longleftrightarrow\quad
 a^{-1}haP=P,\quad(\forall)\,h\in H.
\]
\[
\Longleftrightarrow\quad a^{-1}Ha\leq P
\quad\Longrightarrow\quad H\leq aPa^{-1}.
\]

\medskip
\noindent Let $Q\in\operatorname{Syl}_p(G)$. Then
\[
aQa^{-1}\leq P.
\]
\noindent But
\[
|P|=p^t,\qquad |Q|=p^t
\quad\Longrightarrow\quad |aQa^{-1}|=p^t
\quad\Longrightarrow\quad aQa^{-1}=P.
\]

\bigskip
\noindent\underline{Application}: $|G|=pq$, $p>q$, $q\nmid(p-1)$.

\medskip
\noindent By Cauchy's theorem, there exist $a,b\in G$ with
$o(a)=p$, $o(b)=q$.

\noindent $P=\langle a\rangle$, $Q=\langle b\rangle$.

\noindent By Sylow's Third Theorem:
\[
|\operatorname{Syl}_p(G)|\mid pq,
\qquad |\operatorname{Syl}_p(G)|\equiv1\pmod p
\quad\Longrightarrow\quad |\operatorname{Syl}_p(G)|=1.
\]

\noindent $P\in\operatorname{Syl}_p(G)\ \ni Q$.%
% The final reading "\ni Q" is uncertain. If interpreted as
% containment Q<=P, it conflicts with |Q|=q and |P|=p.
% The manuscript page ends here; no additional proof is supplied.
